% !TEX program = xelatex
\documentclass[11pt,a4paper]{article}

% ---------------------------------------------------------------------------
% Template habituel : note scientifique sobre et argumentee
% ---------------------------------------------------------------------------
\usepackage[a4paper,margin=2.35cm]{geometry}
\usepackage{fontspec}
\usepackage[french]{babel}
\usepackage{microtype}
\usepackage{csquotes}
\usepackage{xcolor}
\usepackage{graphicx}
\usepackage{booktabs}
\usepackage{array}
\usepackage{tabularx}
\usepackage{enumitem}
\usepackage[most]{tcolorbox}
\usepackage{fancyhdr}
\usepackage{titlesec}
\usepackage{setspace}
\usepackage{xurl}
\usepackage{hyperref}
\usepackage{amsmath,amssymb,mathtools,bm}
\usepackage{siunitx}
\usepackage{tikz}
\usetikzlibrary{arrows.meta,positioning,fit,calc,shapes.geometric}
\usepackage{listings}
\setlength{\headheight}{14pt}
\usepackage[backend=biber,style=authoryear,sorting=nyt,maxcitenames=2,maxbibnames=12,giveninits=true]{biblatex}
\addbibresource{references.bib}

\setmainfont{Liberation Serif}
\setsansfont{Liberation Sans}
\setmonofont{Liberation Mono}

\definecolor{mainblue}{HTML}{17324D}
\definecolor{accent}{HTML}{B85C38}
\definecolor{teal}{HTML}{008A8A}
\definecolor{softblue}{HTML}{EEF5FA}
\definecolor{softorange}{HTML}{FFF4EC}
\definecolor{softteal}{HTML}{ECF8F7}
\definecolor{softgray}{HTML}{F6F6F6}
\definecolor{darkgray}{HTML}{333333}

\hypersetup{
  colorlinks=true,
  linkcolor=mainblue,
  citecolor=accent,
  urlcolor=mainblue,
  pdftitle={Cinématique DIC plein champ, calcul local/global et identification par défaut d'équilibre},
  pdfauthor={Jean-François Witz},
  pdfsubject={Plasticité cristalline, DIC, EBSD, local/global, pénalisation, défaut d'équilibre, identification},
  pdfkeywords={DIC, EBSD, plasticité cristalline, local-global, domain decomposition, equilibrium gap, FEMU, data-driven}
}

\pagestyle{fancy}
\fancyhf{}
\lhead{\textsc{DIC plein champ -- stratégie local/global}}
\rhead{\thepage}
\renewcommand{\headrulewidth}{0.4pt}

\titleformat{\section}
  {\Large\bfseries\color{mainblue}}
  {\thesection}{0.75em}{}
\titleformat{\subsection}
  {\large\bfseries\color{darkgray}}
  {\thesubsection}{0.75em}{}
\titleformat{\subsubsection}
  {\normalsize\bfseries\color{darkgray}}
  {\thesubsubsection}{0.75em}{}

\setlist[itemize]{leftmargin=1.3em,itemsep=0.2em,topsep=0.25em}
\setlist[enumerate]{leftmargin=1.5em,itemsep=0.25em,topsep=0.25em}
\setstretch{1.06}
\sisetup{locale=FR,detect-all=true}

\newtcolorbox{thesisbox}[1][]{
  breakable,
  colback=softblue,
  colframe=mainblue,
  boxrule=0.8pt,
  arc=2mm,
  left=1.0em,right=1.0em,top=0.8em,bottom=0.8em,
  #1
}
\newtcolorbox{warningbox}[1][]{
  breakable,
  colback=softorange,
  colframe=accent,
  boxrule=0.8pt,
  arc=2mm,
  left=1.0em,right=1.0em,top=0.8em,bottom=0.8em,
  #1
}
\newtcolorbox{methodbox}[1][]{
  breakable,
  colback=softgray,
  colframe=darkgray,
  boxrule=0.5pt,
  arc=2mm,
  left=1.0em,right=1.0em,top=0.8em,bottom=0.8em,
  #1
}
\newtcolorbox{ideabox}[1][]{
  breakable,
  colback=softteal,
  colframe=teal,
  boxrule=0.8pt,
  arc=2mm,
  left=1.0em,right=1.0em,top=0.8em,bottom=0.8em,
  #1
}

\newcommand{\udic}{\bm{u}^{\mathrm{DIC}}}
\newcommand{\epsdic}{\bm{\varepsilon}^{\mathrm{DIC}}}
\newcommand{\sigmadic}{\bm{\sigma}^{\mathrm{DIC}}}
\newcommand{\jump}[1]{\left[\!\left[ #1 \right]\!\right]}
\newcommand{\Rmech}{\bm{R}}
\newcommand{\Req}{\bm{r}}
\newcommand{\state}{\mathcal{S}}

\begin{document}

\begin{titlepage}
  \centering
  \vspace*{1.2cm}
  {\Huge\bfseries\color{mainblue}
  Exploiter une cinématique DIC plein champ\par}
  \vspace{0.35cm}
  {\LARGE\bfseries
  pour le calcul local/global, le diagnostic de modèle et l'identification constitutive\par}
  \vspace{0.9cm}

  \begin{thesisbox}
  \large
  \textbf{Idée directrice.} Lorsque le déplacement est mesuré dans tout le domaine à chaque instant,
  il ne doit pas être traité comme une condition de Dirichlet forte intérieure. Il peut au contraire
  fournir (i) un prédicteur cinématique et constitutif plein champ extrêmement riche, (ii) une
  régularisation souple temporaire pour découpler et initialiser les sous-domaines, puis (iii) un
  défaut d'équilibre directement exploitable pour évaluer, localiser et éventuellement identifier
  une loi constitutive sans résoudre un problème FEM global complet à chaque essai de paramètres.
  \end{thesisbox}

  \vspace{0.8cm}
  {\large Note d'idéation scientifique et d'architecture numérique\par}
  \vspace{0.25cm}
  {\normalsize Démonstrateur 316L -- DIC / EBSD / plasticité cristalline\par}

  \vfill
  {\large Jean-François Witz\par}
  {\normalsize Univ. Lille, CNRS, Centrale Lille, UMR 9013 -- LaMcube\par}
  \vspace{0.45cm}
  {\normalsize Version du 13 août 2026\par}
\end{titlepage}

\tableofcontents
\newpage

\section*{Résumé exécutif}
\addcontentsline{toc}{section}{Résumé exécutif}

Le point de départ est inhabituel pour un problème de mécanique numérique : la corrélation
 d'images donne une estimation du déplacement à chaque noeud ou pixel du domaine observé et
à plusieurs instants. Dans un calcul direct classique, seule une partie de cette information serait
utilisée comme condition aux limites. L'imposer fortement partout serait pourtant excessif : le
problème mécanique n'aurait pratiquement plus la possibilité de corriger les incompatibilités dues
au bruit de mesure, aux hypothèses 2D, à une loi constitutive imparfaite ou à la discrétisation.

La proposition développée ici est de transformer cette information redondante en un avantage
algorithmique et scientifique. Avant toute décomposition du domaine, le champ DIC complet est
projeté avec la cinématique discrète du solveur, puis la loi matériau est rejouée en chaque point à
partir de l'état convergé au temps précédent. On obtient ainsi un \emph{prédicteur constitutif
expérimental} de l'état interne, de la contrainte, des tangentes et, pour une loi non locale, du champ
micromorphique. Ce prédicteur n'est jamais commité : il sert uniquement d'initialisation.

Le domaine est ensuite partitionné. Les déplacements DIC intérieurs ne deviennent pas des
conditions de Dirichlet : ils apparaissent comme une attache faible pondérée par une raideur
$\kappa_D$. Cette raideur est initialement forte pour tirer chaque sous-domaine vers l'état
expérimental, puis elle décroît lorsque l'équilibre des tractions d'interface exige que la mécanique
s'écarte de la mesure. Dans le mode direct visé ici, $\kappa_D$ doit tendre vers zéro à l'intérieur :
la solution finale ne conserve comme conditions cinématiques imposées que les véritables bords de
Dirichlet. Le champ DIC a alors servi de continuation, de préconditionneur physique et de prédicteur
multiéchelle, sans biaiser la solution finale.

Enfin, le même rejeu constitutif de la cinématique mesurée fournit presque gratuitement un défaut
d'équilibre faible. Intégré dans le temps, ce défaut devient (i) une carte de discordance
expérience--modèle et (ii) une fonction coût d'identification qui ne nécessite pas de résoudre un
problème FEM non linéaire global pour chaque jeu de paramètres. Cette voie se rattache à
l'Equilibrium Gap Method, aux approches d'identification plein champ et à la Constitutive Relation
Error, mais la combinaison proposée avec états internes incrémentaux de plasticité cristalline,
EBSD, non-localité et continuation local/global constitue une piste méthodologique spécifique à
évaluer.

\begin{ideabox}
\textbf{Hypothèse centrale à tester.} Le champ DIC plein champ peut jouer trois rôles successifs :
\emph{prédicteur constitutif global}, \emph{coarse predictor expérimental pour la résolution
local/global}, puis \emph{source d'un défaut d'équilibre utilisé pour le diagnostic et l'inversion}.
Ces trois rôles utilisent la même donnée et la même loi matériau, mais répondent à trois objectifs
numériques différents.
\end{ideabox}

\section{Pourquoi le problème est différent d'une décomposition de domaine classique}

\subsection{Une cinématique expérimentale disponible dans tout le domaine}

À chaque instant $t_n$, on dispose d'un champ expérimental
\begin{equation}
  \udic_n = \udic(\bm{x},t_n)
\end{equation}
sur une grille très dense, potentiellement de plusieurs millions de pixels. La DIC plein champ est
précisément l'une des raisons pour lesquelles les méthodes d'identification mécanique peuvent
exploiter bien davantage d'information que des jauges ou quelques capteurs ponctuels
\parencite{hild2006review}.

La tentation naturelle consiste à imposer $\udic_n$ comme condition de Dirichlet partout. Cela
supprime cependant une grande partie du problème mécanique. Une donnée expérimentale comporte
du bruit et peut être incompatible avec les hypothèses du modèle ; une imposition forte transfère
alors cette incompatibilité directement aux déformations et aux contraintes.

L'objectif est donc différent :
\begin{equation}
  \boxed{
  \text{utiliser la DIC partout pendant la résolution, mais ne conserver in fine
  que les vrais bords de Dirichlet.}
  }
\end{equation}

\subsection{La donnée expérimentale comme information algorithmique et non comme contrainte}

Les approches de régularisation mécanique de la DIC ont déjà montré l'intérêt de combiner
information image et admissibilité mécanique au moyen d'une pénalisation ou d'un filtre mécanique
\parencite{rethore2009mechanicalfilter,mendoza2019regularization}. De même, des travaux ont
explicitement couplé des sous-domaines de DIC régularisée et formulé un problème d'interface
compatible avec le calcul parallèle \parencite{bouclier2017domain}.

La proposition présente inverse toutefois le point de vue. La DIC n'est plus l'inconnue que la
mécanique régularise : elle est déjà mesurée et sert à accélérer un problème direct fortement non
linéaire. La mécanique doit progressivement pouvoir s'en écarter pour atteindre son propre équilibre.

\section{Prédicteur constitutif expérimental avant partitionnement}

\subsection{Projection cinématique cohérente avec le solveur}

Avant de partitionner le domaine, on calcule les déformations expérimentales avec \emph{le même
opérateur discret} que celui du solveur :
\begin{equation}
  \epsdic_n = B_h\,\udic_n .
\end{equation}
Cette étape est préférable à l'import d'un champ de déformations calculé par une chaîne différente :
elle préserve les conventions de cisaillement, les points matériels, les sous-cellules et les filtres
cinématiques du problème numérique.

\subsection{Rejeu incrémental de la loi matériau}

Soit $\state_{n-1}^{\mathrm{conv}}$ l'état matériau commité après convergence mécanique de l'instant
précédent. Pour chaque point matériel $q$, le prédicteur est obtenu indépendamment par
\begin{equation}
  \left(
  \state_{n,q}^{\mathrm{pred}},
  \bm{\sigma}_{n,q}^{\mathrm{pred}},
  \mathbb{C}_{n,q}^{\mathrm{pred}},
  \Gamma_{n,q}^{\mathrm{pred}}
  \right)
  =
  \mathcal{M}_{\bm{\theta}}
  \left(
  \state_{n-1,q}^{\mathrm{conv}},
  \epsdic_{n,q},
  \Delta t_n
  \right),
  \label{eq:predictor}
\end{equation}
où $\bm{\theta}$ désigne les paramètres constitutifs.

Pour une plasticité cristalline, $\state$ contient notamment les glissements et variables de
durcissement ; $\Gamma$ peut représenter une mesure scalaire d'activité plastique, par exemple la
somme des glissements cumulés positifs. Le calcul de \eqref{eq:predictor} est massivement parallèle
par point matériel.

\begin{warningbox}
\textbf{Règle de transaction impérative.} L'état $\state_n^{\mathrm{pred}}$ n'est jamais commité.
Pendant la résolution mécanique, toute réévaluation doit repartir de
$\state_{n-1}^{\mathrm{conv}}$ et intégrer directement jusqu'à la cinématique trial courante.
Commiter le prédicteur DIC puis le corriger créerait un faux morceau de chemin de chargement en
plasticité path-dependent.
\end{warningbox}

\subsection{Extension non locale}

Si la loi dépend d'un champ micromorphique $\chi$, le prédicteur peut inclure un petit point fixe
\emph{sans résoudre l'équilibre mécanique} :
\begin{align}
  \Gamma_n^{(k)} &= \mathcal{G}\left(\state_{n-1}^{\mathrm{conv}},\epsdic_n,\chi_n^{(k)}\right),\\
  (I-\ell^2\Delta)\chi_n^{(k+1)} &= \Gamma_n^{(k)}.
\end{align}
La cinématique étant figée par la DIC, le coût reste constitué d'intégrations matérielles pointwise et
d'un problème de Helmholtz scalaire. Le résultat fournit $\chi_n^{\mathrm{pred}}$ avant le début du
problème local/global.

\subsection{Ce que le prédicteur fournit gratuitement}

Avant toute itération d'interface, on possède déjà une approximation cohérente de
\begin{equation}
\left\{
\udic_n,
\epsdic_n,
\state_n^{\mathrm{pred}},
\bm{\sigma}_n^{\mathrm{pred}},
\mathbb{C}_n^{\mathrm{pred}},
\Gamma_n^{\mathrm{pred}},
\chi_n^{\mathrm{pred}}
\right\}.
\end{equation}
En particulier, après partitionnement, chaque interface hérite immédiatement de tractions
prédites $\bm{t}^{\mathrm{pred}}=\bm{\sigma}^{\mathrm{pred}}\bm{n}$. Le défaut initial de continuité des
tractions est donc mesurable sans solve global.

\begin{figure}[htbp]
\centering
\begin{tikzpicture}[
  node distance=8mm and 12mm,
  box/.style={draw=mainblue, rounded corners=2mm, fill=softblue, align=center, minimum height=9mm, text width=32mm},
  mat/.style={draw=teal, rounded corners=2mm, fill=softteal, align=center, minimum height=9mm, text width=32mm},
  glob/.style={draw=accent, rounded corners=2mm, fill=softorange, align=center, minimum height=9mm, text width=34mm},
  arr/.style={-{Latex[length=2.2mm]}, thick, draw=darkgray}
]
\node[box] (dic) {DIC plein champ\\$\udic_n(\bm{x})$};
\node[box, right=of dic] (kin) {cinématique discrète\\$\epsdic_n=B_h\udic_n$};
\node[mat, right=of kin] (umat) {rejeu UMat/MFront\\depuis $\state_{n-1}^{\rm conv}$};
\node[mat, below=of umat] (pred) {prédicteur global\\$\state_n^{pred},\bm{\sigma}_n^{pred},\chi_n^{pred}$};
\node[glob, left=of pred] (part) {partitionnement\\restrictions par sous-domaine};
\node[glob, left=of part] (lg) {local/global\\pénalisation DIC décroissante};
\node[box, below=of lg] (conv) {solution convergée\\Dirichlet seulement au bord};
\draw[arr] (dic)--(kin);
\draw[arr] (kin)--(umat);
\draw[arr] (umat)--(pred);
\draw[arr] (pred)--(part);
\draw[arr] (part)--(lg);
\draw[arr] (lg)--(conv);
\end{tikzpicture}
\caption{Ordre logique proposé. Le prédicteur constitutif est construit sur le champ global avant que la partition ne définisse des interfaces artificielles. Les états prédits sont ensuite distribués aux sous-domaines comme initialisation, sans commit constitutif.}
\label{fig:pipeline}
\end{figure}

\section{Résolution local/global avec pénalisation DIC décroissante}

\subsection{Problème local pénalisé}

Après partitionnement $\Omega=\cup_s\Omega_s$, chaque sous-domaine résout un problème non
linéaire incluant une attache faible au champ expérimental :
\begin{equation}
  \Rmech_s(\bm{u}_s,\state_s)
  + \kappa_D^{(k)}\,\bm{W}_{D,s}
    \left(\bm{u}_s-\udic_s\right)
  + \bm{R}_{\Gamma,s} = \bm{0}.
  \label{eq:penalizedlocal}
\end{equation}
$\bm{R}_{\Gamma,s}$ représente le couplage d'interface et $\bm{W}_{D,s}$ pondère la confiance
accordée aux données. Dans le cas idéal, $\bm{W}_D$ est relié à l'inverse de la covariance de mesure
plutôt qu'à un poids uniforme ; cette logique rejoint l'importance de la métrique statistique dans
l'identification plein champ \parencite{roux2020optimal}.

Dans une implémentation matrix-free, la contribution au Jacobien est immédiate :
\begin{equation}
  \bm{J}_s\bm{v}
  = \bm{J}_{\mathrm{mech},s}\bm{v}
  + \kappa_D^{(k)}\bm{W}_{D,s}\bm{v}
  + \bm{J}_{\Gamma,s}\bm{v}.
\end{equation}
Elle ne nécessite aucune matrice globale assemblée.

\subsection{Deux contraintes d'interface distinctes}

Sur une interface $\Gamma_{ss'}$, la solution mécanique doit tendre vers
\begin{align}
  \jump{\bm{u}} &= \bm{u}_s-\bm{u}_{s'} = \bm{0},\\
  \jump{\bm{t}} &= \bm{t}_s+\bm{t}_{s'} = \bm{0}.
\end{align}
La pénalisation DIC ne doit pas être confondue avec le mécanisme de couplage des interfaces. Elle
apporte une information extérieure au problème mécanique ; l'interface impose l'admissibilité
globale.

\subsection{Continuation ``data vers mécanique''}

L'idée centrale est d'utiliser une raideur de pénalisation variable :
\begin{equation}
  \kappa_D^{(0)} > \kappa_D^{(1)} > \cdots > \kappa_D^{(K)}.
\end{equation}
Au début, $\kappa_D$ est suffisamment grand pour maintenir les solutions locales près de la
cinématique expérimentale. Les sous-domaines sont donc fortement initialisés et presque découplés.
Lorsque le saut de traction ne peut plus diminuer sans s'écarter du champ mesuré, la pénalisation est
relâchée.

Pour le problème direct visé, la cible est
\begin{equation}
  \boxed{\kappa_{D,\mathrm{int}}^{(K)}\rightarrow 0}
\end{equation}
sur les degrés de liberté intérieurs. La solution finale n'est alors pas une solution assimilée : elle
est une solution mécanique dont les seules conditions cinématiques permanentes sont celles des
bords physiques de Dirichlet.

\begin{methodbox}
\textbf{Alternative assimilation.} Un second mode pourrait conserver un plancher
$\kappa_D^{\mathrm{stat}}>0$ fixé par le niveau de bruit expérimental. Ce mode répondrait à une autre
question : trouver l'état mécaniquement admissible le plus compatible avec la mesure compte tenu de
son incertitude. Il faut le distinguer explicitement du problème direct non biaisé.
\end{methodbox}

\subsection{Règle adaptative de relâchement}

Trois indicateurs peuvent piloter la continuation :
\begin{align}
 \eta_D &= \left[\frac{1}{N_D}(\bm{u}-\udic)^T\bm{C}_D^{-1}(\bm{u}-\udic)\right]^{1/2},\\
 \eta_u &= \frac{\|\jump{\bm{u}}\|}{u_{\mathrm{ref}}},\\
 \eta_t &= \frac{\|\jump{\bm{t}}\|}{t_{\mathrm{ref}}}.
\end{align}
Une règle de départ pourrait être : si $\eta_D$ est inférieur ou comparable au niveau de bruit mais
que $\eta_t$ stagne au-dessus de sa cible, diminuer $\kappa_D$. Autrement dit, on ne demande pas au
calcul de reproduire la mesure plus précisément que ce que sa propre incertitude justifie lorsque
cela empêche l'équilibre mécanique.

Le code du projet contient déjà une calibration de pénalité au bord par principe de discrépance : la
raideur est fixée de façon que l'écart RMS simulation--DIC atteigne l'incertitude expérimentale, et
non pour améliorer artificiellement la convergence. La présente stratégie généralise cette idée à
l'intérieur du domaine, mais avec un rôle temporaire de continuation.

\subsection{Le champ DIC comme coarse predictor expérimental}

Dans une décomposition de domaine classique, les modes de grande longueur d'onde motivent un
espace grossier afin que l'information traverse rapidement le domaine. Ici, $\udic$ contient déjà une
estimation globale des translations, rotations, extensions moyennes et grandes variations spatiales.
Il peut donc être interprété comme un \emph{coarse predictor expérimental distribué}. Cette idée ne
supprime pas a priori la nécessité d'un problème grossier, mais elle conduit à une hypothèse testable :
le nombre d'itérations globales pourrait devenir beaucoup moins sensible au nombre de
sous-domaines.

Les méthodes LATIN multiscalaires montrent précisément l'intérêt d'un niveau macro pour accélérer
la propagation d'information entre sous-structures \parencite{ladeveze2010latin}. Une contribution
potentielle du présent cadre serait d'étudier dans quelle mesure la cinématique mesurée peut
compléter ou réduire cet espace macro.

\subsection{Algorithme proposé à un instant de mesure}

\begin{methodbox}
\textbf{Algorithme conceptuel au temps $t_n$.}
\begin{enumerate}
  \item Charger $\udic_n$ et calculer $\epsdic_n=B_h\udic_n$ sur le domaine complet.
  \item Depuis $\state_{n-1}^{\mathrm{conv}}$, rejouer la loi matériau à tous les points pour obtenir
  $\state_n^{\mathrm{pred}}$, $\bm{\sigma}_n^{\mathrm{pred}}$ et, si nécessaire, $\chi_n^{\mathrm{pred}}$.
  \item Partitionner le domaine et distribuer aux sous-domaines les restrictions du prédicteur.
  \item Initialiser les inconnues locales avec $\udic_n$ et un $\kappa_D$ élevé.
  \item Résoudre les sous-domaines en parallèle sans commiter les états matériaux.
  \item Calculer les sauts de déplacement et de traction aux interfaces ; effectuer la correction globale.
  \item Si l'équilibre d'interface stagne alors que l'écart à la DIC est inférieur au niveau statistiquement utile,
  diminuer $\kappa_D$ et recommencer.
  \item Lorsque les résidus locaux et les interfaces satisfont les critères et que la pénalisation intérieure a
  atteint sa cible (zéro pour le problème direct), commiter une seule fois $\state_n^{\mathrm{conv}}$.
\end{enumerate}
\end{methodbox}

\section{Le prédicteur DIC comme détecteur presque gratuit d'insuffisance du modèle}

\subsection{État constitutivement admissible, mais pas nécessairement statiquement admissible}

Le rejeu de la loi sur $\epsdic$ produit un état compatible avec la loi choisie, mais rien ne garantit
que le champ de contraintes correspondant satisfasse l'équilibre. Pour les degrés de liberté
intérieurs, on peut calculer le résidu faible
\begin{equation}
  \Req_n(\bm{\theta})
  = \bm{f}_{\mathrm{int}}
  \left(\bm{\sigma}^{\mathrm{DIC}}_n(\bm{\theta})\right)
  - \bm{f}_{\mathrm{ext},n}.
  \label{eq:eqgap}
\end{equation}
En forme élémentaire,
\begin{equation}
  \bm{r}_{e,n}
  = \int_{\Omega_e} \bm{B}^T\bm{\sigma}^{\mathrm{DIC}}_n\,\mathrm{d}\Omega
  - \bm{f}_{e,n}^{\mathrm{ext}}.
\end{equation}
Il est préférable d'utiliser cette forme faible FEM plutôt que de calculer directement
$\nabla\!\cdot\!\bm{\sigma}$, qui amplifierait les dérivations du bruit expérimental.

Cette logique est directement apparentée à l'Equilibrium Gap Method, introduite pour identifier des
champs d'endommagement à partir de mesures de déplacement plein champ
\parencite{claire2004equilibrium}. Des extensions ont ensuite exploité les résidus et sensibilités
pour diagnostiquer la qualité d'un modèle et guider son enrichissement
\parencite{neggers2017enrichment}.

\subsection{Carte spatio-temporelle de discordance expérience--modèle}

Pour ne pas juger une loi sur un seul instant, on peut accumuler l'indicateur dans le temps :
\begin{equation}
  E_e(\bm{\theta})
  = \sum_{n=1}^{N_t} w_n\,
    \bm{r}_{e,n}^T\bm{W}_{r,e,n}\bm{r}_{e,n}.
  \label{eq:localcum}
\end{equation}
La carte $E_e$ indique les zones où, sur toute l'histoire, la cinématique observée conduit la loi à
produire un champ de contraintes difficilement compatible avec l'équilibre.

\begin{warningbox}
Un grand $E_e$ n'est pas une preuve que la seule loi constitutive est erronée. Le défaut peut aussi
provenir du bruit DIC, d'un mauvais recalage DIC--EBSD, d'effets hors-plan, d'une épaisseur ou de
forces extérieures mal connues, de la discrétisation ou d'une hypothèse cinématique. La quantité doit
être interprétée comme une \textbf{discordance expérience--modèle} avant toute attribution causale.
\end{warningbox}

La pondération $\bm{W}_r$ peut être reliée à la covariance du résidu induite par l'incertitude DIC.
Cette normalisation permettrait de distinguer une discordance mécaniquement significative d'un
résidu compatible avec le bruit de mesure, dans l'esprit des métriques statistiques proposées pour
l'identification plein champ \parencite{roux2020optimal}.

\subsection{Lecture physique avec EBSD}

Dans le démonstrateur 316L, la carte $E_e$ peut être croisée avec :
\begin{itemize}
  \item l'orientation cristalline et la désorientation ;
  \item les joints de grains ;
  \item l'activité des systèmes de glissement ;
  \item le glissement cumulé $\Gamma$ ;
  \item les gradients de $\Gamma$ ou du champ micromorphique $\chi$ ;
  \item les zones de localisation observées par DIC.
\end{itemize}
Une concentration du défaut aux joints de grains, dans certains grains ou dans les bandes de
localisation devient alors une information directe pour comparer J2, SRIX local, SRIX non local ou
d'autres enrichissements constitutifs.

\section{Ouverture inverse : identification par écart à l'équilibre}

\subsection{Pourquoi cette voie peut être beaucoup moins chère que la FEMU}

Une FEMU classique évalue chaque jeu de paramètres $\bm{\theta}$ par
\begin{equation}
 \bm{\theta}
 \longrightarrow
 \text{solve FEM non linéaire complet}
 \longrightarrow
 \bm{u}^{\mathrm{calc}}(\bm{\theta})
 \longrightarrow
 \|\bm{u}^{\mathrm{calc}}-\udic\|.
\end{equation}
L'identification intégrée et la FEMU sont des outils puissants, y compris pour les lois
élastoplastiques \parencite{bertin2016idic}, mais elles impliquent généralement des résolutions
mécaniques répétées.

Ici, la DIC fournit déjà le champ cinématique. Pour une loi candidate, l'évaluation devient
\begin{equation}
 \bm{\theta}
 \longrightarrow
 \text{rejeu constitutif local de l'histoire DIC}
 \longrightarrow
 \bm{\sigma}^{\mathrm{DIC}}(\bm{\theta})
 \longrightarrow
 \Req(\bm{\theta}).
\end{equation}
La fonction coût peut être
\begin{equation}
  J_{\mathrm{eq}}(\bm{\theta})
  = \sum_{n=1}^{N_t}
  \Req_n(\bm{\theta})^T
  \bm{C}_{r,n}^{-1}
  \Req_n(\bm{\theta}).
  \label{eq:objective}
\end{equation}
Pour une loi locale, le calcul dominant est alors la chronologie constitutive en chaque point ; tous
les points sont parallèles spatialement. Pour la loi non locale, il faut ajouter le ou les solves de
Helmholtz sur $\chi$, mais on évite toujours le Newton/Krylov mécanique global pour chaque jeu de
paramètres.

\begin{table}[htbp]
\centering
\caption{Comparaison conceptuelle entre FEMU et pré-identification par défaut d'équilibre.}
\begin{tabularx}{\textwidth}{@{}>{\bfseries}p{3.2cm}XX@{}}
\toprule
 & FEMU / I-DIC mécanique & Rejeu DIC + défaut d'équilibre \\
\midrule
Cinématique & calculée par le modèle & imposée par la mesure pour l'évaluation \\
Coût par $\bm{\theta}$ & solve mécanique global non linéaire & intégration constitutive + assemblage du résidu \\
Parallélisme & limité par le solve global & massif par point matériel \\
État interne & issu du solve direct & rejoué incrémentalement sur l'histoire DIC \\
Sortie locale & écart déplacement / image & défaut d'équilibre local spatio-temporel \\
Rôle recommandé & validation finale / raffinement & filtrage, initialisation, exploration paramétrique \\
\bottomrule
\end{tabularx}
\end{table}

Il ne faut pas supposer a priori que $J_{\mathrm{eq}}$ conduit exactement au même optimum que la
FEMU. La stratégie la plus robuste est donc hiérarchique :
\begin{equation}
  \boxed{
  \text{identification rapide par }J_{\mathrm{eq}}
  \;\longrightarrow\;
  \text{quelques candidats}
  \;\longrightarrow\;
  \text{validation par solve mécanique complet.}
  }
\end{equation}

\subsection{Sensibilités et optimisation}

Si l'on dispose des sensibilités constitutives
$\partial\bm{\sigma}/\partial\bm{\theta}$, alors
\begin{equation}
  \frac{\partial J_{\mathrm{eq}}}{\partial\theta_i}
  = 2\sum_n
  \Req_n^T\bm{C}_{r,n}^{-1}
  \frac{\partial \Req_n}{\partial\theta_i},
\end{equation}
avec
\begin{equation}
  \frac{\partial \Req_n}{\partial\theta_i}
  = \mathcal{A}
  \left(
  \frac{\partial\bm{\sigma}^{\mathrm{DIC}}_n}{\partial\theta_i}
  \right),
\end{equation}
où $\mathcal{A}$ est l'opérateur faible d'équilibre. L'essentiel de la sensibilité reste donc local au
matériau et n'exige pas nécessairement un adjoint mécanique global. Cette piste peut devenir
particulièrement attractive avec un backend GPU différentiable ou des sensibilités fournies par la
loi matériau.

\begin{figure}[htbp]
\centering
\begin{tikzpicture}[
  node distance=9mm and 12mm,
  box/.style={draw=mainblue, rounded corners=2mm, fill=softblue, align=center, minimum height=9mm, text width=34mm},
  eq/.style={draw=accent, rounded corners=2mm, fill=softorange, align=center, minimum height=9mm, text width=35mm},
  outbox/.style={draw=teal, rounded corners=2mm, fill=softteal, align=center, minimum height=9mm, text width=36mm},
  arr/.style={-{Latex[length=2.2mm]}, thick, draw=darkgray}
]
\node[box] (theta) {paramètres candidats\\$\bm{\theta}$};
\node[box, right=of theta] (replay) {rejeu de toute l'histoire DIC\\pointwise et path-dependent};
\node[eq, right=of replay] (res) {résidu faible d'équilibre\\$\Req_n(\bm{\theta})$};
\node[outbox, below=of res] (cost) {coût global\\$J_{eq}(\bm{\theta})$};
\node[outbox, left=of cost] (map) {carte locale cumulée\\$E_e(\bm{\theta})$};
\node[box, left=of map] (femu) {quelques candidats\\validation FEMU / solve complet};
\draw[arr] (theta)--(replay);
\draw[arr] (replay)--(res);
\draw[arr] (res)--(cost);
\draw[arr] (res)--(map);
\draw[arr] (cost)--(femu);
\draw[arr] (map)--(femu);
\end{tikzpicture}
\caption{Ouverture inverse. La cinématique mesurée remplace le solve mécanique comme source des gradients de déformation pendant la phase d'exploration paramétrique. Le solve complet reste l'étape de validation finale.}
\label{fig:inverse}
\end{figure}

\section{Positionnement scientifique et hypothèse de nouveauté}

La proposition ne part pas d'une page blanche. Plusieurs familles de travaux couvrent des morceaux
du problème :
\begin{itemize}
  \item l'Equilibrium Gap Method exploite l'équilibre et les champs de déplacement mesurés pour
  identifier des paramètres ou champs constitutifs \parencite{claire2004equilibrium};
  \item la Constitutive Compatibility Method a été formulée dans une décomposition de domaine
  spécifiquement destinée aux mesures plein champ, avec problèmes locaux de type Dirichlet puis
  restauration de l'équilibre macroscopique \parencite{lubineau2015dd};
  \item la DIC régularisée mécaniquement et son couplage de domaines ont déjà été formulés avec des
  inconnues d'interface et une perspective parallèle \parencite{bouclier2017domain,mendoza2019regularization};
  \item la Constitutive Relation Error fournit un cadre récent de modélisation data-driven où les
  relations physiques et les données expérimentales sont combinées explicitement
  \parencite{ladeveze2024cre};
  \item les champs résiduels peuvent servir à diagnostiquer l'insuffisance d'un modèle et à guider
  des enrichissements progressifs \parencite{neggers2017enrichment}.
\end{itemize}

\begin{ideabox}
\textbf{Hypothèse de contribution originale, à confirmer par une revue bibliographique exhaustive.}
La piste spécifique n'est pas l'usage de la DIC, de l'equilibrium gap, de la pénalisation ou de la
décomposition de domaine pris séparément. Elle réside dans leur assemblage pour un problème direct
et inverse path-dependent :
\begin{enumerate}
  \item prédire avant partitionnement les états internes d'une loi non linéaire à partir de l'histoire DIC ;
  \item utiliser cette prédiction comme initialisation globale cohérente des sous-domaines ;
  \item appliquer une attache DIC intérieure temporaire, relâchée en fonction de l'équilibre d'interface jusqu'à disparition dans le problème direct ;
  \item réutiliser exactement le même rejeu constitutif pour construire une fonction coût d'équilibre et une carte de défaut de modèle ;
  \item appliquer le tout à une plasticité cristalline hétérogène informée par EBSD, potentiellement non locale et massivement parallèle.
\end{enumerate}
\end{ideabox}

La question méthodologique peut alors être formulée ainsi :
\begin{quote}
\emph{Une cinématique expérimentale dense peut-elle remplacer une partie du rôle du coarse space
numérique dans une résolution local/global non linéaire, puis servir de support à une identification
constitutive par défaut d'équilibre sans solve mécanique répété ?}
\end{quote}

\section{Architecture numérique compatible avec le plein EBSD}

Le plein champ EBSD de plusieurs milliers de pixels par direction dépasse l'échelle des crops de
validation. L'architecture proposée est néanmoins favorable au calcul distribué :
\begin{itemize}
  \item le prédicteur constitutif DIC est massivement parallèle par point matériel ;
  \item les sous-domaines peuvent être résolus indépendamment pendant les phases locales ;
  \item le solveur matrix-free évite le stockage d'une matrice globale ;
  \item seules les traces d'interface, les corrections globales et éventuellement le champ $\chi$ doivent circuler entre sous-domaines ;
  \item la pénalisation DIC est une opération diagonale, naturelle sur GPU.
\end{itemize}

Une implémentation Dask/CUDA peut donc être envisagée à grain grossier : un worker par GPU conserve
les états d'un ou plusieurs sous-domaines, tandis que les kernels locaux utilisent CuPy/cuFFT. Le
scheduler ne doit pas intervenir dans chaque produit Jacobien--vecteur ; il orchestre des solves
locaux complets et des étapes globales. Le prédicteur DIC peut être calculé avant le partitionnement
au sens algorithmique, puis distribué par blocs sans modifier sa définition globale.

\section{Programme de validation proposé}

\subsection{Étape 1 -- preuve de concept J2 sur un domaine où la solution complète est connue}

Utiliser un cas M100 ou M200 pour lequel le solve monolithique actuel fournit une référence. Découper
artificiellement le domaine en $2\times2$, $4\times4$, puis davantage.

Comparer :
\begin{enumerate}
  \item solve local/global sans information DIC intérieure ;
  \item simple initialisation $\bm{u}^0=\udic$ ;
  \item prédicteur constitutif DIC sans pénalisation ;
  \item pénalisation DIC fixe ;
  \item pénalisation DIC décroissante jusqu'à zéro ;
  \item éventuellement pénalisation décroissante + coarse space classique.
\end{enumerate}

Mesurer le nombre d'itérations globales, le nombre total de Newton/Krylov locaux, le coût matériau,
les sauts de traction/déplacement, l'écart final à la référence et la sensibilité au partitionnement.

\subsection{Étape 2 -- validation de l'indicateur d'équilibre}

Pour la même histoire DIC :
\begin{enumerate}
  \item rejouer la loi sans solve global et calculer $\Req_n$ ;
  \item vérifier que le résidu du champ issu du solve convergé est nettement plus faible ;
  \item comparer les cartes cumulées pour plusieurs lois ou paramètres connus ;
  \item injecter des erreurs contrôlées dans les paramètres et vérifier que $J_{\mathrm{eq}}$ possède un minimum près des paramètres de référence ;
  \item propager un bruit DIC synthétique pour définir $\bm{C}_r$ et les seuils de significativité.
\end{enumerate}

\subsection{Étape 3 -- SRIX et EBSD}

Une fois le principe qualifié en J2, l'intérêt scientifique apparaît avec la plasticité cristalline :
rejeu path-dependent des systèmes de glissement, croisement des cartes de résidu avec l'EBSD et
comparaison local/non-local. Le prédicteur expérimental peut également être utilisé pour fournir les
états initiaux des futures résolutions distribuées sur la carte EBSD complète.

\subsection{Étape 4 -- inversion rapide puis validation FEMU}

Sur un petit nombre de paramètres SRIX :
\begin{enumerate}
  \item cartographier $J_{\mathrm{eq}}(\bm{\theta})$ ;
  \item lancer une optimisation sans solve mécanique global ;
  \item retenir quelques candidats ;
  \item exécuter seulement ensuite le solve complet pour comparer les champs calculés à la DIC et les réactions aux mesures de force disponibles.
\end{enumerate}

Le rapport de coût entre les deux étages devra être mesuré et non supposé. Un gain important est
attendu conceptuellement, mais sa magnitude dépendra du coût du matériau, du nombre de temps DIC,
des éventuelles itérations non locales et du nombre de paramètres.

\section{Questions ouvertes à conserver}

\begin{methodbox}
\begin{enumerate}
  \item Quelle loi de décroissance de $\kappa_D$ maximise la convergence sans introduire un nouveau paramètre arbitraire ?
  \item Peut-on relier directement la décision de relâchement au principe de discrépance et au rapport $\eta_t/\eta_D$ ?
  \item La DIC permet-elle réellement de réduire la dimension du coarse space, ou seulement d'améliorer l'initialisation ?
  \item Faut-il pondérer $\bm{W}_D$ point par point avec une covariance DIC estimée, ou une covariance simplifiée suffit-elle ?
  \item Comment séparer dans $E_e$ le défaut constitutif d'un défaut de recalage DIC--EBSD ou d'un effet 3D ?
  \item Pour SRIX non local, doit-on conserver un solve Helmholtz global pendant les premières versions ou décomposer également $\chi$ ?
  \item Jusqu'où $J_{\mathrm{eq}}$ conserve-t-il l'identifiabilité de paramètres qui seraient observables en FEMU ?
  \item Les sensibilités constitutives suffisent-elles à obtenir un gradient robuste de $J_{\mathrm{eq}}$ pour une optimisation de grande dimension ?
\end{enumerate}
\end{methodbox}

\section{Conclusion}

La disponibilité d'un champ DIC dense à chaque instant change la structure même du problème. Le
champ expérimental ne doit pas être réduit à un chargement de bord, ni imposé fortement dans le
volume. Il peut d'abord être transformé en un prédicteur complet de la cinématique et des états
internes, puis servir de guide temporaire à une résolution local/global qui libère progressivement
les degrés de liberté intérieurs jusqu'à retrouver un problème mécanique classique aux seules
conditions de bord.

Cette même opération de rejeu constitutif possède une seconde valeur : elle produit sans solve
global un défaut d'équilibre faible, calculable localement, cumulable dans le temps et comparable à
l'incertitude expérimentale. Elle ouvre ainsi une voie de diagnostic spatial de l'insuffisance du
modèle et une stratégie d'identification potentiellement beaucoup moins coûteuse que la FEMU pour
l'exploration paramétrique.

\begin{thesisbox}
\textbf{Vision unifiée.} La DIC plein champ pourrait devenir simultanément :
\begin{center}
\textbf{un prédicteur d'état} $\rightarrow$ \textbf{un accélérateur local/global}
$\rightarrow$ \textbf{un détecteur de défaut de modèle}
$\rightarrow$ \textbf{une fonction coût d'identification}.
\end{center}
Le caractère convaincant de cette proposition dépend maintenant de trois démonstrations :
absence de biais du solve direct lorsque la pénalisation intérieure disparaît, réduction effective du
coût de résolution avec le nombre de sous-domaines, et pouvoir discriminant du défaut d'équilibre
pour les paramètres et enrichissements constitutifs.
\end{thesisbox}

\printbibliography

\end{document}
